% OLD TEXT FROM HERE ON
% This chapter contains text from the original submission.
% It will be revised as part of the salvage plan.

\section{Data and exploratory empirical setup} \label{sec:data}

To see and compare how competing model variants based on IO-tables differ in terms of outcomes, one obvious approach is to impose some exogenous shock and compare the reactions of different models to these shocks. However, given the different paradigmatic perspectives associated with different model variants, it is not too surprising to observe, that the archetypical understanding of what, exactly, constitutes an exogenous shock also differs across traditions. From the perspective of a classic Leontief model the most obvious form of an exogenous shock is given by a shift in final demand (which, in turn, leads to restructuring of the supply side). While the Leontief framework can, in principle, also be applied to simulate the effects of other types of shocks -- such as shocks to prices (e.g., by using the Leontief price equation, i.e. equation \eqref{eq:supply_clear}) or technological shocks (e.g. by manipulating the technical coefficients collected in $A$), the most widely used application focuses on simulating the impacts of a change in final demand. Such an approach is obviously based on a specific "output-adjustment" model closure \citep{Taylor1990}, that is characteristic for typical applications of the model, not only because of the Keynesian features emphasized in section \ref{subsec:Leontief}, but also because this type of analysis is suitable to cover basic scenarios for policy-analysis evaluating different types of public investments or prospective changes in public spending.

Similarly, the archetypical understanding of an exogenous shock in a more supply-side driven neoclassical vision is best conceived as a sudden change in aggregate productivity as captured by $\mathbf{\mathcal A}$. This approach is most intuitive as production is seen as inherently constraining final demand, which is, in principle, conceptualized as open-ended as more is always deemed better in such a framework. The technology parameter $\mathbf{\mathcal A}$ affects these crucial production conditions in an exhaustive way and, thereby, emerges as a preferred entry point for exogenous shocks in neoclassical models. Nonetheless, the basic structure of the model is, again, able to accommodate a greater variety of shocks, including exogenous changes to final demand, in principle, if suitable additional conditions and parameters are introduced.

These aspects are essential as in our application the type of shock is derived from the core research interest, namely to assess how paradigmatic differences translate into different estimates on the economic impacts associated with additional investments undertaken to confront climate warming in particular and ecological degradation in general. Hence, we start from a shift in final demand as this adequately represents our conceptualization of an exogenous shock, when asking for the economic impacts of ecologically motivated investments.

To provide a concrete and reliable empirical example for such investment requirements, we take the German residential housing sector as our main case of analysis. When doing so we base our analysis and discussion on a specific estimate of the costs associated with a transformation of the German housing sector towards carbon neutrality as supplied by \citet{Hornykewycz.2025}. This study estimates expected costs by proportionally scaling past renovations costs under the assumption that an increase in the renovation intensity takes place, that is sufficient to meet Germany's climate goals. Specifically, the study at hand finds that renovation efforts will have to be more than doubled, if Germany is set to comply with its overall goal to achieve carbon neutrality till 2050. Such a shift would amount to yearly costs of about 40.3 bn. € in 2019 prices\footnote{We assumed a deflator of 1.46 based on price indices for construction costs as specified in \citet{Hornykewycz.2025}. In 2023 prices the shock amounts to about 58 bn. €.}, which would represent a substantial shift in final demand. As the study also details how these costs should be mapped on different sectors, it provides sufficient information to construct a concrete and directly applicable specification for the sector-specific annual changes in final demand imposed by such a reorientation. This specification shocks final demand in seven sectors, where sectors are sorted by the relative increase associated with these shocks (as given in brackets): 
Glass and glassware (+50.3\%), ceramic products and processed stones (+37,3 \%), specialized construction work (+27,8\%), rubber and plastic products (+13.4\%), chemicals and chemical products (+3.7\%), machinery (+1.2\%) and electrical equipment (+1.0\%).
In terms of the data used, we apply the so estimated shock vector as proposed by \citet{Hornykewycz.2025} with official German input-output data for 2019, the most recent year available \citep[see][]{destatis_strukturdaten}.

In this study we exploit the opportunity to draw on a well-researched suggestions for how a specific transformation strategy impacts on final demands across sectors, which perfectly resembles our envisioned scenario of an exogenous shock to final demands. As both modeling approaches introduced can cope with such a simulated change in final demands, this provides a clear-cut case for illustrating \textit{what difference the choice of modeling philosophy makes} when approaching practical issues related to the ambition to speed up socio-ecological transformation by additional public or private investment.

% HIER KÖNNTEN WIR BERNI FRAGEN, OB ER WAS SCHREIBEN MAG. WIR KÖNNTEN AUCH ÜBERLEGEN, DIE SCHIENENAPPLICATION MIT AUFZUNEHMEN.

